\section{Related Work}
\label{sec:related}

\subsection{Motor Activity as a Psychiatric Biomarker}

The use of actigraphy to characterize psychiatric conditions has a well-established empirical foundation. Garcia-Ceja et al.\ \cite{garcia2018depresjon} introduced the Depresjon dataset that forms the basis of our work, alongside a suite of baseline classifiers that demonstrated the feasibility of distinguishing mood-disorder patients from healthy controls using wrist-worn activity counts. Their best-performing approaches relied on random forests applied to hand-engineered statistical descriptors extracted over fixed windows. Zanella-Calzada et al.\ \cite{zanella2019feature} extended this direction by developing a systematic pipeline for feature extraction from motor activity signals, extracting over twenty statistical and frequency-domain features and demonstrating improved classification performance over the original baselines. Jakobsen et al.\ \cite{jakobsen2020applying} subsequently applied a broader suite of machine learning methods to the same dataset, with the explicit goal of discriminating between bipolar and unipolar patients and healthy controls. Their work highlighted the particular difficulty of the bipolar-versus-unipolar sub-task and noted that classifiers performed markedly worse when the target was differential diagnosis rather than simple healthy-versus-depressed discrimination. All three of these works operate within the classical machine learning paradigm: a domain expert first defines a feature set, and a classifier is then fit to those features. Our work departs from this tradition by training a deep neural network end-to-end on raw one-minute activity counts, removing domain-driven feature selection entirely.

\subsection{Deep Learning for Wearable Time-Series Classification}

The application of deep learning to wearable sensor data has a rich recent history in the Human Activity Recognition (HAR) literature. LeCun et al.\ \cite{lecun2015deep} established the general principle that deep architectures can learn hierarchical representations from raw sensor data more effectively than hand-crafted features in many domains. Yang et al.\ \cite{yang2015deep} demonstrated that 1D convolutional networks applied directly to raw accelerometer streams matched or exceeded hand-engineered feature pipelines on standard HAR benchmarks. Ord\'{o}{\~n}ez and Roggen \cite{ordonez2016deep} extended this finding by showing that stacking convolutional layers on top of LSTM modules produces robust activity classifiers that generalize across different sensor placements and body-worn configurations. Fawaz et al.\ \cite{fawaz2019deep} conducted a comprehensive review of deep learning approaches for time-series classification, confirming that CNN-LSTM hybrids are among the strongest architectures for wearable sensing tasks. A key distinction between HAR tasks and the psychiatric classification problem we address is temporal scale: HAR models classify discrete physical behaviors occurring over seconds, while we aim to infer a clinical diagnosis from movement evolution over 10 to 14 days. This extended temporal horizon is what motivates the LSTM component in our architecture, as it provides gated memory capable of integrating information across multiple days of recording \cite{hochreiter1997lstm}.

\subsection{Attention Mechanisms and Bidirectional Sequence Models}

Beyond the standard unidirectional LSTM, several architectural extensions have been proposed to improve temporal representation learning. Graves and Schmidhuber \cite{graves2005bilstm} introduced bidirectional recurrent networks for sequence modeling, demonstrating that processing sequences in both forward and backward directions simultaneously allows each time step to be represented in the context of both its past and its future. Bahdanau et al.\ \cite{bahdanau2015attention} introduced the attention mechanism as a way for neural networks to selectively focus on the most relevant portions of an input sequence when making a prediction, rather than relying on a single compressed hidden state. Vaswani et al.\ \cite{vaswani2017attention} generalized this idea into the Transformer architecture, demonstrating that attention alone can match or surpass recurrent models on sequential tasks. In the context of psychiatric actigraphy, attention could in principle allow the model to identify diagnostically relevant time-of-day windows, such as sleep-onset timing or early-morning activity, that carry disproportionate information about diagnosis. We evaluate both bidirectional processing and attention-augmented architectures in Approach~2 of our experimental investigation.

\subsection{Circadian Rhythm Disruption and Mood Disorders}

A substantial body of psychiatric literature links disruptions to the sleep-wake cycle to both bipolar and unipolar depression, and argues that the two conditions produce distinct circadian signatures. Wehr et al.\ \cite{wehr1987sleep} documented phase advances in circadian rhythms as a characteristic feature of bipolar states, with bipolar patients showing earlier sleep onset and a compressed sleep-wake cycle compared to controls. Boudebesse et al.\ \cite{boudebesse2014sleep} demonstrated day-by-day correlations between rest-activity timing variability and mood state in bipolar patients in remission, suggesting that actigraphy captures real-time fluctuations in mood-linked physiological processes. Goodwin and Jamison \cite{goodwin2007manic} provide a comprehensive account of the phenotypic overlap between bipolar depression and unipolar depression, noting that the depressive phase of bipolar disorder is behaviorally indistinguishable from unipolar depression without collateral history of manic or hypomanic episodes. This literature provides the primary theoretical motivation for the LSTM component of our architecture, and also motivates our investigation of 48-hour and 72-hour temporal windows in Approach~3, since single 24-hour windows may be too short to capture the multi-day rhythm shifts characteristic of bipolar disorder.

\subsection{Class Imbalance in Clinical Machine Learning}

Severe class imbalance is a pervasive challenge in clinical classification tasks, where the condition of interest is often rarer than the control class. Chawla et al.\ \cite{chawla2002smote} proposed the Synthetic Minority Oversampling Technique (SMOTE), which generates synthetic minority-class examples by interpolating between existing real examples in feature space, as a general-purpose remedy for this problem. SMOTE has since become a standard component of clinical machine learning pipelines. However, He and Garcia \cite{he2009learning} raise important caveats: if the minority-class examples are not a reliable sample from the true data distribution, as may be the case when the class has only eight members as in our bipolar group, synthetic interpolation can create artificial examples that do not correspond to real physiological patterns and may degrade generalization. Our Baseline approach evaluates SMOTE in exactly this setting, and our Approach~1 investigates whether downsampling the majority class to achieve a natural 1:1 balance is a more principled alternative. These two experiments together allow us to decouple the effects of class imbalance from underlying signal quality.
